%---------------------------------------
% Codificación e idioma
%---------------------------------------
\usepackage[T1]{fontenc}
% \usepackage[utf8]{inputenc} % Innecesario en compilaodres modernos (LaTeX post-2018 / Tectonic), UTF-8 es nativo.
\usepackage[spanish,es-tabla]{babel}

%---------------------------------------
% Márgenes del documento
%---------------------------------------
\usepackage[
  letterpaper,
  left=3cm,        % 3 cm a la izquierda para encuadernación
  right=2.5cm,     % 2.5 cm a la derecha
  top=2.5cm,       % 2.5 cm arriba
  bottom=2.5cm     % 2.5 cm abajo
]{geometry}

%---------------------------------------
% Fuente Times New Roman
%---------------------------------------
\usepackage{newtxtext,newtxmath}

%---------------------------------------
% Espaciado y párrafos
%---------------------------------------
\usepackage{setspace}
\onehalfspacing % O \doublespacing si la universidad te exige interlineado doble estricto

\usepackage{indentfirst}
\setlength{\parindent}{1.27cm} % Sangría de primera línea
\setlength{\parskip}{0pt}      % Sin espacio extra entre párrafos (estándar académico)

%---------------------------------------
% Tipografía y Alineación
%---------------------------------------
\usepackage{microtype}
\usepackage{ragged2e}

\AtBeginDocument{
  \justifying
}

%---------------------------------------
% Bibliografía
%---------------------------------------

\usepackage[
  backend=biber,
  style=apa
]{biblatex}

% Cargamos el archivo .bib (reemplaza 'database.bib' por el nombre exacto de tu archivo)
\addbibresource{database.bib}
%---------------------------------------
% Tablas
%---------------------------------------
\usepackage{array}
\usepackage{booktabs}
\usepackage{multirow}
\usepackage{longtable}
\usepackage{caption}

%---------------------------------------
% Figuras
%---------------------------------------
\usepackage{graphicx}
\usepackage{float}
\usepackage{subcaption} % 'subcaption' reemplaza al paquete obsoleto 'subfigure'
\usepackage{epstopdf}

%---------------------------------------
% Matemáticas y Encabezados
%---------------------------------------
\usepackage{bm}
\usepackage{fancyhdr}
\pagestyle{fancy}
\fancyhf{} % Limpia todos los encabezados y pies de página actuales

% Quitamos la línea horizontal del encabezado
\renewcommand{\headrulewidth}{0pt} 

% Colocamos el número de página al centro del pie de página
\rhead{\thepage}
%---------------------------------------
% Formato de títulos
%---------------------------------------
\usepackage{titlesec}

\titleformat{\chapter}[hang]
  {\centering\normalfont\fontsize{12pt}{14pt}\selectfont\bfseries} % \centering al inicio
  {\thechapter.}              % Etiqueta (ej: 1., 2.)
  {0.5em}                     % Espacio entre el número y el título
  {}

  % 2. Formato para Secciones (\section) a 12pt
\titleformat{\section}
  {\normalfont\fontsize{12pt}{14pt}\selectfont\bfseries}
  {\thesection}
  {1em}
  {}

% 3. Formato para Subsecciones (\subsection) a 12pt
\titleformat{\subsection}
  {\normalfont\fontsize{12pt}{14pt}\selectfont\bfseries}
  {\thesubsection}
  {1em}
  {}

% 2. Controlamos los espacios alrededor del título
% SINTAXIS: \titlespacing*{comando}{izquierda}{espacio_arriba}{espacio_abajo}
\titlespacing*{\chapter}{0pt}{-1.5cm}{15pt}
\titlespacing*{\section}{0pt}{12pt}{6pt}
\titlespacing*{\subsection}{0pt}{10pt}{4pt}

%---------------------------------------
% Hipervínculos
%---------------------------------------
\usepackage[
  colorlinks=true,
  linkcolor=blue,
  citecolor=black,
  urlcolor=blue
]{hyperref}

%---------------------------------------
% Otros paquetes
%---------------------------------------
\usepackage{multicol}
\usepackage{titling}
\usepackage{anyfontsize}